\documentclass[11pt]{article}
\usepackage[T1]{fontenc}
\usepackage{textcomp}
\usepackage[margin=0.75in]{geometry}
\usepackage{amsmath,amssymb,amsthm}
\usepackage{array,booktabs}
\usepackage{hyperref}
\setlength{\parindent}{0pt}
\newtheorem{theorem}{Theorem}
\theoremstyle{definition}
\newtheorem{definition}{Definition}
\title{Flow Proof Artifact: taylor-sin-maclaurin}
\author{Generated by \texttt{flow doc proof}}
\date{Flow Proof Book}
\begin{document}
\maketitle
Near the origin, sin(x) agrees with each successive Maclaurin partial sum.\\[0.5em]
\textbf{Source.} brook taylor — https://en.wikipedia.org/wiki/Taylor\_series\\[0.5em]
\bigskip
\noindent{\large \textbf{Derived fact 1.} \textit{Near the origin, sin(x) agrees with each successive Maclaurin partial sum} \hfill real analysis \cdot Taylor series \cdot sin equals its Maclaurin series near zero}\par
\smallskip\noindent\textit{Coordinate: real analysis \cdot Taylor series \cdot sin equals its Maclaurin series near zero}\par
\smallskip\noindent\textit{Source: brook taylor — https://en.wikipedia.org/wiki/Taylor\_series}\par
\smallskip\noindent\textit{Built on: the derivatives of sine at zero follow the alternating pattern of the Maclaurin series}\par
\medskip
\noindent\textbf{Goal.} Near the origin, sin(x) agrees with each successive Maclaurin partial sum\par
\noindent$\displaystyle \forall x \in \mathbb{R}\quad S_{5}(x) = \sin(x) \quad (x \to 0)$\par
\medskip
\noindent
\begin{tabular}{@{} >{\raggedright\arraybackslash}p{0.06\textwidth} >{\raggedright\arraybackslash}p{0.47\textwidth} >{\raggedleft\arraybackslash}p{0.41\textwidth} @{}}
\textbf{\#} & \textbf{Proof} & \textbf{Mathematics} \\
\midrule
\textbf{1.} & We prove that sin equals its Maclaurin series near zero for Taylor series on real analysis. & \\[0.45em]
\textbf{2.} & We invoke the definitional clause governing smooth functions on real analysis: the derivatives of sine at zero follow the alternating pattern of the Maclaurin series. & \\[0.45em]
\textbf{3.} & From step 2, the Maclaurin coefficients match the known derivatives of sine at zero, so each partial sum agrees with sin(x) to the next order, so taylor sin(x, order 5) equals sin(x) in a neighbourhood of zero. Hence proven. & $\displaystyle S_{5}(x) = \sin(x) \quad (x \to 0)$ \\[0.45em]
\bottomrule
\end{tabular}
\label{thm:Analysis-Taylor series-sin_equals_its_maclaurin_series_near_zero}
\par\medskip\hrule
\end{document}
